Large language model (LLM) agents are increasingly evaluated on tasks whose value depends on \emph{verifiable} numerical outcomes. Mathematical modeling---the practice of turning an open-ended real-world problem into assumptions, a model, a solution, and a defended recommendation---is a natural but previously unmeasured target: it is the format of the CUMCM and MCM contests, it is multimodal (official statements ship figures and spreadsheets), and its answers are only meaningful if they are \emph{both} close to a credible reference \emph{and} physically feasible.

We introduce \bench{}, a benchmark of 18 real mathematical-modeling problems (CUMCM 2023--2025 and MCM 2023--2025, problems A/B/C each year). Each task delivers the verbatim contest statement, the official attachments, and the official figures cropped from the source PDFs; the agent must produce one structured JSON artifact, scored only on the numeric core of the \emph{final} contest question against a reproduction of the corresponding outstanding (O-award) paper.

Two design choices distinguish \bench{}. First, scoring is \emph{direction-aware} and \emph{one-sided}: a metric declares whether more is better, less is better, or a target is sought, so a model is never penalised for beating the paper, while degenerate scoring bands and non-comparable reference units are handled by explicit, auditable rules. Second, feasibility is verified twice---by a programmatic gate covering all 18 tasks through six decidable invariant forms, and by a Lean~4 certificate layer that discharges 624 obligations on exact integers with no Mathlib, no \texttt{sorry}, and no (kernel-trusting) \texttt{native\_decide}. A hard gate zeroes answers that violate a stated physical constraint.

We evaluate 16 model/harness configurations, including five multimodal agents that can read the official figures. The strongest entry reaches \oeval{} 68.07\% (codex + GPT-5.5), while the cheapest near-frontier entry reaches 65.76\% at \$2.54 per 18-task run. We report per-task behaviour, a cost/robustness analysis, and three case studies (a top proprietary model, the O-award reproduction, and a low-cost open model) on representative problems. We further document---and, where the fault was ours, fix---the ways a reproduction-anchored benchmark can silently penalise correct answers, and we show how the dual feasibility layer makes the remaining failures machine-checkable rather than merely asserted.
